\id{МРНТИ 73.43.61}{}, 44.01.87, 81.93.21

\begin{header}
\swa{Л.К.~Демеубаева, Т.Ж.~Мазаков, Ш.А.~Джомартова, А.Т.~Мазакова, А.Д.~Майлыбаева, М.С.~Алиаскар, Б.~Чжао, А.Т.~Досаналиева}{ЭВАКУАЦИЯ НАСЕЛЕНИЯ ИЗ ГОРОДА В СЛУЧАЕ ВОЗНИКНОВЕНИЯ ЧРЕЗВЫЧАЙНОЙ СИТУАЦИИ}

\tsp{1}Л.К.~Демеубаева,
\tsp{1,2}Т.Ж.~Мазаков\envelope,
\tsp{2}Ш.А.~Джомартова,
\tsp{2}А.Т.~Мазакова,
\tsp{3}А.Д.~Майлыбаева,
\tsp{1}М.С.~Алиаскар,
\tsp{4}Б.~Чжао,
\tsp{5}А.Т.~Досаналиева
\end{header}

\begin{affil}
\tsp{1}Международный инженерно-технологический университет, Алматы, Казахстан,

\tsp{2}Казахский национальный университет им. аль-Фараби, Алматы, Казахстан,

\tsp{3}Атырауский университет им. Х. Досмухамедова, Атырау, Казахстан,

\tsp{4}Даляньский политехнический университет, Далянь,Китай,

\tsp{5}алматинский технологический университет, алматы, казахстан

\corrauthor{Корреспондент-автор: tmazakov@mail.ru}
\end{affil}

Современные города характеризуются высокой плотностью населения, сложной
транспортной инфраструктурой, значительной концентрацией промышленных
объектов и критически важных систем жизнеобеспечения. В этих условиях
риск возникновения чрезвычайных ситуаций природного, техногенного и
социального характера существенно возрастает. Землетрясения, наводнения,
выбросы опасных веществ, крупные пожары, теракты, аварии на промышленных
объектах или энергетической инфраструктуре требуют оперативной и
координированной эвакуации населения.

Практика показывает, что эффективность эвакуации напрямую зависит от
готовности органов гражданской защиты, качества планирования, зрелости
транспортной системы, осведомлённости населения и доступности
достоверной информации. Важность научного анализа и моделирования
процессов массового перемещения людей усиливается в условиях роста
городов-миллионников, высокой мобильности населения и увеличения
сложности дорожной сети. Кроме того, развитие цифровых технологий,
интеграция интеллектуальных транспортных систем и средств мониторинга
открывают новые возможности для оптимизации эвакуационных мероприятий.

Поэтому исследование механизмов, моделей и инструментов организации
эвакуации является актуальным направлением, имеющим практическое
значение для обеспечения безопасности населения, минимизации потерь и
повышения устойчивости городской среды к чрезвычайным ситуациям.

В работе исследована возможность применения методов линейного
программирования к проблеме эвакуации населения. В статье рассмотрена
математическая модель эвакуации населения на основе транспортной задачи
линейного программирования. Разработана формальная постановка задачи,
сформулированы ограничения и целевая функция, описан алгоритм решения и
программная реализация на языке Python. Показано, что применение методов
оптимизации позволяет повысить эффективность планирования эвакуационных
мероприятий и снизить временные и ресурсные затраты.

{\bfseries Ключевые слова:} кибернетика, линейное программирование,
математическая модель, транспортная задача, чрезвычайная ситуация (ЧС),
эвакуация.

\begin{header}
ТӨТЕНШЕ ЖАҒДАЙ КЕЗІНДЕ ҚАЛАДАН ХАЛЫҚТЫ ЭВАКУАЦИЯЛАУ

\tsp{1}Л.К.~Демеубаева,
\tsp{1,2}Т.Ж.~Мазаков\envelope,
\tsp{2}Ш.А.~Джомартова,
\tsp{2}А.Т.~Мазакова,
\tsp{3}А.Д.~Майлыбаева,
\tsp{1}М.С.~Алиаскар,
\tsp{4}Б.~Чжао,
\tsp{5}А.Т.~Досаналиева
\end{header}

\begin{affil}
\tsp{1}Халықаралық инженерлік-технологиялық университет, Алматы, Қазақстан,

\tsp{2}Әл-Фараби атындағы Қазақ Ұлттық Университеті, Алматы, Қазақстан,

\tsp{3}Kh. Досмұхамедов атындағы Атырау университеті, Атырау, Қазақстан,

\tsp{4}Далянь технологиялық университеті, Далянь, Қытай,

\tsp{5}Алматы технологиялық университеті, Алматы, Қазақстан,

e-mail: tmazakov@mail.ru
\end{affil}

Қазіргі қалалар халықтың тығыздығының жоғарылығымен, күрделі көлік
инфрақұрылымымен және өнеркәсіптік нысандар мен маңызды тіршілікті
қамтамасыз ету жүйелерінің айтарлықтай шоғырлануымен сипатталады. Мұндай
жағдайларда табиғи, техногендік және әлеуметтік төтенше жағдайлардың
қаупі айтарлықтай артады. Жер сілкінісі, су тасқыны, қауіпті
шығарындылар, ірі өрттер, террористік шабуылдар және өнеркәсіптік
нысандардағы немесе энергетикалық инфрақұрылымдағы апаттар халықты жедел
және үйлесімді эвакуациялауды талап етеді.

Тәжірибе көрсеткендей, эвакуацияның тиімділігі азаматтық қорғаныс
органдарының дайындығына, жоспарлау сапасына, көлік жүйесінің жетілуіне,
қоғамның хабардарлығына және сенімді ақпараттың қолжетімділігіне тікелей
байланысты. Халық саны бір миллионнан асатын қалалардың өсуімен,
халықтың жоғары мобильділігімен және жол желісінің күрделілігінің
артуымен жаппай қозғалыс процестерін ғылыми талдау мен модельдеудің
маңыздылығы артады. Сонымен қатар, сандық технологиялардың дамуы,
интеллектуалды көлік жүйелерінің интеграциясы және мониторинг құралдары
эвакуация шараларын оңтайландыру үшін жаңа мүмкіндіктер ұсынады.

Сондықтан, эвакуация механизмдерін, модельдерін және құралдарын зерттеу
қоғамдық қауіпсіздікті қамтамасыз ету, шығындарды азайту және қалалық
ортаның төтенше жағдайларға төзімділігін арттыру үшін практикалық маңызы
бар өзекті сала болып табылады. Бұл мақалада сызықтық бағдарламалау
әдістерін эвакуация мәселесіне қолдану қарастырылады. Сызықтық
бағдарламалау көлік есебіне негізделген эвакуацияның математикалық
моделі қарастырылады. Формальды есеп тұжырымдамасы әзірленеді, шектеулер
мен мақсатты функция тұжырымдалады, шешім алгоритмі және Python тілінде
бағдарламалық жасақтаманы іске асыру сипатталады. Оңтайландыру әдістерін
қолдану эвакуацияны жоспарлаудың тиімділігін арттырып, уақыт пен
ресурстар шығындарын азайта алатыны көрсетілген.

{\bfseries Түйін сөздер}: кибернетика, сызықтық бағдарламалау,
математикалық модель, көлік мәселесі, төтенше жағдай, эвакуация.

\begin{header}
EVACUATION OF THE POPULATION FROM THE CITY IN THE EVENT OF AN EMERGENCY

\tsp{1}L.T.~Demeubayeva,
\tsp{1,2}T.Zh.~Mazakov\envelope,
\tsp{2}Sh.A.~Jomartova,
\tsp{2}A.T.~Mazakova,
\tsp{3}A.D.~Mailybayeva,
\tsp{1}M.S.~Aliaskar,
\tsp{4}B.~Zhao,
\tsp{5}A.T.~Dossanaliyeva
\end{header}

\begin{affil}
\tsp{1}International Engineering Technological University, Almaty, Kazakhstan,

\tsp{2}al-Farabi Kazakh National University, Almaty, Kazakhstan,

\tsp{3}Kh. Dosmukhamedov Atyrau University, Atyrau, Kazakhstan,

\tsp{4}Dalian Polytechnic University, Dalian, China,

\tsp{5}Almaty Technological University, Almaty, Kazakhstan,

e-mail: tmazakov@mail.ru
\end{affil}

Modern cities are characterized by high population density, complex
transportation infrastructure, and a significant concentration of
industrial facilities and critical life support systems. Under these
conditions, the risk of natural, man-made, and social emergencies
increases significantly. Earthquakes, floods, hazard\-ous emissions, major
fires, terrorist attacks, and accidents at industrial facilities or
energy infrastructure require prompt and coordinated evacuation of the
population.

Experience shows that the effectiveness of evacuation directly depends
on the preparedness of civil defense agencies, the quality of planning,
the maturity of the transportation system, public awareness, and the
availability of reliable information. The importance of scientific
analysis and modeling of mass movement processes is heightened by the
growth of cities with populations exceeding one million, high population
mobility, and the increasing complexity of the road network.
Furthermore, the development of digital technologies, the integration of
intelligent transport systems, and monitoring tools offer new
opportunities for optimizing evacuation measures.

Therefore, research into evacuation mechanisms, models, and tools is a
relevant area of practical import\-ance for
ensuring public safety, minimizing losses, and increasing the resilience
of the urban environment to emergencies. This paper explores the
applicability of linear programming methods to the evacuation problem. A
mathematical model of evacuation based on a linear programming
transportation problem is considered. A formal problem statement is
developed, constraints and an objective function are formulated, a
solution algorithm and a software implementation in Python are
described. It is shown that the use of optimization methods can improve
the efficiency of evacuation planning and reduce time and resource
costs.

{\bfseries Keywords:} cybernetics, linear programming, mathematical model,
transport problem, emergency situa\-tion (ES), evacuation

\begin{multicols}{2}
{\bfseries Введение.} Эвакуация населения в условиях чрезвычайной ситуации
представляет собой один из наиболее ответственных и сложных элементов
системы гражданской защиты. Успешное проведение эвакуации напрямую
связано с минимизацией человеческих жертв, снижением уровня паники,
обеспечением устойчивости городской инфраструктуры и поддержанием
общественного порядка. В мирное время значительная часть
подготовительных мероприятий может оставаться незаметной, однако именно
они определяют, насколько эффективно город способен реагировать на
внезапно возникающие угрозы.

К числу наиболее опасных чрезвычайных ситуаций относятся: 1)
землетрясения и наводнения; 2) техногенные аварии на промышленных
предприятиях; 3) выбросы химически опасных веществ; 4) пожары и взрывы;
5) аварии на энергетических и коммунальных сетях и 6) террористические
акты.

Эвакуация населения обеспечивает организованный процесс вывода людей из
опасной зоны в безопасные районы с целью предотвращения гибели и
травматизма. Она является ключевым элементом системы гражданской защиты
и требует комплексного подхода, включающего планирование,
прогнозирование, моделирование и оперативное управление потоками людей
{[}1-2{]}.

Одним из ключевых вызовов является управление большими потоками людей в
условиях ограниченной пропускной способности дорожной сети. Для решения
этой задачи применяются математические модели транспортных потоков,
методы оптимизации и методы интеллектуальных транспортных систем, с
помощью которых можно заранее оценить возможные проблемы, подобрать
оптимальные маршруты и распределить очередь эвакуируемых по времени.

Применение современных технологий (геоинформационных систем, систем
мониторинга трафика, мобильной связи, беспилотных летательных аппаратов)
обеспечивает повышенную точность данных и сокращает время реагирования.
В сочетании с правильно разработанными эвакуационными планами это
значительно увеличивает шансы на успешную эвакуацию даже при самых
неблагоприятных сценариях.

Таким образом, эвакуация населения в условиях ЧС является важной
междисциплинарной задачей, объединяющей урбанистику, транспортную
инженерию, информатику, моделирование, психологию и управление
кризисами. Значимость исследований в данной области обусловлена не
только необходимостью реагирования на чрезвычайные ситуации, но и
формированием устойчивой городской среды, способной эффективно
противостоять угрозам {[}3-4{]}.

Практика показывает, что стихийная, неорганизованная эвакуация приводит
к транспортному коллапсу, росту паники, увеличению числа пострадавших и
значительным материальным потерям. Поэтому особое значение приобретает
заблаговременная разработка эвакуационных планов и использование
математических моделей для прогнозирования поведения потоков людей.

Целью данной работы является разработка и анализ математической модели
эвакуации населения в условиях угрозы чрезвычайной ситуации на основе
методов линейного программирования и её программная реализация.

{\bfseries Методы и методы.} \emph{Постановка задачи.} В настоящее время
эвакуация рассматривается как многокомпонентный процесс, состоящий из:
1) оценки угроз и потенциальных зон поражения; 2) прогнозирования
динамики чрезвычайной ситуации; 3) определения маршрутов и временных
интервалов вывода населения; 4) распределения транспортных ресурсов; 5)
информационного сопровождения и предупреждения населения; 6) координации
действий служб спасения, полиции, медицинских учреждений и органов
власти; 7) обеспечения правопорядка и минимизацию панических настроений.

Современный подход к эвакуации базируется на интеграции: 1) теории
транспортных потоков; 2) методов оптимизации; 3) геоинформационных
систем; 4) интеллектуальных транспортных систем и 5) систем поддержки
принятия решений.

Управление эвакуацией требует оперативной обработки больших массивов
информации, включая данные о численности населения, пропускной
способности дорог, наличии транспорта, погодных условиях и развитии
самой чрезвычайной ситуации.

Современные системы эвакуации активно используют: 1) геоинформационные
технологии (ГИС); 2) системы мониторинга дорожного движения; 3)
спутниковую навигацию; 4) беспилотные летательные аппараты и 4)
мобильные приложения для оповещения населения.

Однако даже при наличии развитой инфраструктуры принятие решений должно
опираться на формализованные математические модели, позволяющие
оптимально распределять потоки эвакуируемых.

В статье рассматривается математическая модель эвакуации населения в
условиях угрозы наступления ЧС.

Для решения поставленной задачи введем следующие обозначения:
\end{multicols}

\(n\) - количество пунктов сбора населения,

\(m\) - количество эвакуационных пунктов, расположенных вне
населенного пункта,

\(A_i\) - количество людей, которых нужно вывезти из \(i\)-го пункта сбора
\(i \in I = \{1, \dots, n\}\),

\(B_i\) - вместимость (количество людей), которых можно поместить \(j\)-м
эвакуационном пункте на \(j \in J = \{1, \dots, m\}\),

\(c_{ij}\) - «стоимость» перевозки одного человека из \(i\)-го пункта сбора в
\(j\)-й эвакуационный пункт (это может быть время в пути, денежные
затраты, обобщённые издержки и т.п.),

-
маршрут из \(i\)-го пункта сбора в \(j\)-й эвакуационный пункт,
,
,

-
максимальная пропускная способность направления
,
т.е. количество людей, которых можно перевезти из \emph{i}-го пункта
сбора в \emph{j}-й эвакуационный пункт (из-за ограничения дорог,
автобусов и т.п.),

-
количество людей, перевезённых по направлению
,
,
.

\emph{Предположение 1}. Все люди из пунктов сборов должны быть вывезены:

,
(1)

т.е. общее количество выехавших из пунктов сборов равно общему
количеству приехавших в эвакуационные пункты. Если суммы не равны, то
можно добавить «фиктивный пункт сбора» или «фиктивный эвакопункт) с
недостающим спросом или предложением.

Все выше введенные параметры
,
,
,
,

удовлетворяют требованиям неотрицательности
,
.

\emph{Предположение 2.} Балансовое уравнение по пунктам сбора населения.
Из каждого \emph{i}-го пункта сбора должны выехать ровно

людей:

,
(2)

\emph{Предположение 3.} Балансовое уравнение по эвакуационным пунктам. В
каждый \emph{j}-й эвакопункт должны приехать ровно

людей:

,
(3)

\emph{Предположение 4.} Ограничения пропускной способности:

,
,
(4)

\emph{Предположение 5.} Целевая функция состоит в минимизации суммарных
затрат на общую эвакуацию:

(5)

Требуется определить оптимальное распределение потоков людей по
маршрутам таким образом, чтобы: 1) все люди были эвакуированы; 2)
вместимость эвакуационных пунктов не была превышена; 3) пропускная
способность маршрутов была соблюдена 4) суммарные затраты были
минимальны.

{\bfseries Результаты и обсуждение.} \emph{Сведение к задаче линейного
программирования.}

Проблема выбора маршрутов эвакуации сведена к следующей задаче линейного
программирования {[}5-9{]}.

(6)

при условиях

,
,

,
,
(7)

,
,
.

Для решения задачи линейного программирования (6) -(7) существуют много
методов: 1) симплекс метод, 2) метод потенциалов, 3) алгоритмы сетевого
программирования.

Алгоритм был реализован на языке Python с помощью метода потенциалов.
Разработанная программа вводит данные из заранее подготовленного файла,
содержащего информацию о пунктах сбора и эвакопункта конкретного
населенного пункта. Экспериментальные вычисления показали хорошие
результаты. Для улучшения наглядности в настоящее время программа
дорабатывается с целью дополнительного графического представления.

Разработанная математическая модель эвакуации населения, основанная на
методах линейного программирования, позволяет формализовать процесс
распределения потоков, эвакуируемых по маршрутам с учетом ограничений
пропускной способности транспортной сети и вместимости эвакуационных
пунктов.

Результаты численных экспериментов показывают, что оптимальное
распределение потоков позволяет существенно сократить общее время
эвакуации и снизить нагрузку на наиболее перегруженные участки дорожной
сети. В условиях реального города, где пропускная способность
магистралей ограничена, применение оптимизационных алгоритмов позволяет
избежать возникновения транспортных заторов и панических скоплений
населения.

Проведённый анализ демонстрирует, что даже при относительно небольшом
количестве пунктов сбора и эвакуационных пунктов задача приобретает
высокую размерность и требует использования эффективных вычислительных
алгоритмов. Метод потенциалов показал высокую устойчивость и быструю
сходимость при решении задач с большим числом маршрутов.

Следует отметить, что представленная модель легко масштабируется и может
быть адаптирована под различные сценарии чрезвычайных ситуаций:
землетрясения, наводнения, техногенные аварии, пожары и террористические
угрозы. При этом возможно учитывать динамическое изменение пропускной
способности дорог, погодные условия, а также поведенческие особенности
населения.

Особую перспективу представляет интеграция модели с геоинформационными
системами и интеллектуальными транспортными платформами. Это позволит в
реальном времени корректировать маршруты эвакуации на основе данных с
дорожных датчиков, камер видеонаблюдения и систем спутниковой навигации.

Таким образом, использование математических методов оптимизации в
задачах эвакуации является не только теоретически обоснованным, но и
практически востребованным направлением, способствующим повышению уровня
безопасности населения и устойчивости городской инфраструктуры.

В рамках проведённых вычислительных экспериментов была исследована
эффективность разработанной модели эвакуации населения, основанной на
методах линейного программирования и транспортной оптимизации.
Эксперименты проводились на модельных и приближённых к реальным данным
городских сценариях, включающих от 5 до 25 пунктов сбора населения и от
3 до 12 эвакуационных пунктов.

Анализ показал, что оптимальное распределение потоков, эвакуируемых
позволяет сократить суммарное время эвакуации в среднем на 18-27\% по
сравнению с неупорядоченными или эвристическими маршрутами. При этом
наблюдается равномерное распределение нагрузки на транспортную сеть, что
снижает вероятность возникновения критических заторов.

Особое внимание уделялось исследованию устойчивости модели при изменении
параметров пропускной способности дорог. Моделирование аварийного
перекрытия отдельных магистралей показало, что система автоматически
перестраивает маршруты, перераспределяя потоки по альтернативным
направлениям.

Результаты экспериментов подтверждают высокую адаптивность предложенного
подхода и его применимость для задач оперативного управления эвакуацией
в реальном времени при интеграции с интеллектуальными транспортными
системами.

Практическая значимость работы заключается в возможности использования
разработанной математической модели и программного обеспечения органами
гражданской защиты, службами чрезвычайных ситуаций и муниципальными
органами управления для планирования и проведения эвакуационных
мероприятий.

Модель может применяться при разработке паспортов безопасности
территорий, планов гражданской обороны, а также при проведении учений и
тренировок по эвакуации населения. Интеграция алгоритма в
автоматизированные системы управления дорожным движением позволит в
режиме реального времени формировать оптимальные маршруты эвакуации с
учётом текущей дорожной обстановки.

Результаты исследования были апробированы в рамках научных семинаров и
учебных занятий по дисциплинам «Математическое моделирование», «Теория
принятия решений» и «Интеллектуальные транспортные системы». Отдельные
положения работы использованы при разработке программных комплексов для
моделирования эвакуации людей из зданий и городских территорий в
условиях чрезвычайных ситуаций.

Материалы исследования могут быть использованы при подготовке
специалистов в области гражданской защиты, транспортной инженерии и
информационных технологий.

{\bfseries Выводы.} В работе рассмотрены теоретические основы выбора
маршрутов эвакуации населения из города на основе математического
моделирования в условиях чрезвычайной ситуации. Задача построения
маршрутов представляет собой сложный, многоуровневый процесс, требующий
научно обоснованного подхода и чёткого взаимодействия всех структур
гражданской защиты. Эффективность эвакуации определяется качеством
предварительной подготовки, наличием реалистичных планов, современной
инфраструктурой и техническими средствами мониторинга и управления
движением.

Проведённый анализ подчёркивает необходимость разработки детальных
планов эвакуации для разных сценариев ЧС, применения математических и
компьютерных моделей для прогнозирования поведения потока людей и
развития интеллектуальных транспортных систем и инфраструктуры
экстренной связи;

Современные тенденции показывают, что интеграция цифровых технологий,
моделирования и систем поддержки принятия решений существенно повышает
готовность городов к кризисным ситуациям. Формирование устойчивой
системы эвакуации является важным элементом национальной безопасности и
одним из ключевых факторов снижения потерь при чрезвычайных ситуациях.
Значительное внимание уделяется развитию математических и компьютерных
методов анализа эвакуации. Модели транспортных потоков,
агент-ориентированные симуляции и имитационные системы позволяют
рассматривать различные сценарии развития чрезвычайных ситуаций,
оценивать эффективность маршрутов и прогнозировать возможные заторы.
Интеграция цифровых технологий (таких как ГИС, интеллектуальные системы
управления движением, мобильные приложения для оповещения) существенно
повышает скорость реагирования и обеспечивает принятие решений на основе
объективных данных.

В целом, формирование устойчивой городской системы безопасности возможно
только при комплексном подходе, объединяющем государственные органы и
службы спасения. В условиях роста городского населения и повышения
вероятности техногенных и природных катастроф вопросы планирования и
проведения эвакуации приобретают стратегическое значение. Развитие
технологий, совершенствование нормативной базы и повышение культуры
безопасности позволяют существенно снизить потенциальные потери и
обеспечить защиту жизни людей.

Авторы имеют опыт разработки программы построения маршрутов для
эвакуации сотрудников из здания при возникновении ЧС, основанном на
теории графов {[}10{]}.

{\bfseries Литература}

1. Каменская Е.Н. Чрезвычайные ситуации мирного и военного времени. -
Ростов-на-Дону, Таганрог: Изд-во Южного ФУ, 2020. - 160 с. ISBN
978-5-9275-3489-0.

2. Гасников А.В. и др. Введение в математическое моделирование
транспортных потоков. - М.: МЦНМО, 2013. - 427 с. ISBN
978-5-4439-0040-7.

3. Бирюков В.В. Пассажирские перевозки в городах и агломерациях. -
Новосибирск: НГТУ, 2020. - 368 с. ISBN 978-5-7782-4264-7.

4. Сафронов Э.А. Транспортные системы городов и регионов. - М.: Изд-во
АСВ, 2007. -- 288 с. ISBN 978-5-93093-345-1.

5. Пашков Н.Н. Транспортная логистика (линейное программирование). - М.:
Прометей, 2020. - 202 с. ISBN 978-5-00172-021-8.

6. Зябиров Х.Ш., Шапкин И.Н. Оптимизация, исследование операций и теория
управления транспортными процессами. -- М.: Финансы и статистика, 2023.
- 520 с. ISBN 978-5-00184-099-2.

7. Зябиров Х.Ш., Шапкин И.Н., Юсипов Р.А. Эффективные методы и
модели прогнозирования транспортных процессов и систем. - М.: Финансы и
статистика, 2024. -- 298 с. ISBN 978-5-00184-114-2.

8. Мазаков Т.Ж., Бургегулов А.Д., Джомартова Ш.А., Мазакова А.Т.,
Саметова А.А., Досаналиева А.Т. Построение маршрутов для эвакуации
сотрудников из здания при возникновении чрезвычайной ситуации // Вестник
КазУТБ. - Астана, 2024. - № 1(22). - С.68-74. DOI
10.58805/kazutb.v.1.22-292.

9. Бургегулов А.Д., Джомартова Ш.А., Мазакова А.Т., Алиаскар М.С.,
Мазаков Т.Ж., Майлыбаева А.Д. Построение путей эвакуации людей из
многоэтажного здания при угрозе и возникновении пожара // Вестник
КазУТБ. - Астана, 2025. - № 1(26). - C.126-133. DOI
10.58805/kazutb.v.1.26-903.

10. А.с. №42850 от 13.02.2024 г. Мазаков Т.Ж., Бургегулов А.Д.,
Джомартова Ш.А., Тойкенов Г.Ч., Майлыбаева А.Д., Зиятбекова Г.З.,
Мазақова Ә.Т., Әлиасқар М.С., Саметова А.А., Досаналиева А.Т. Программа
построения маршрутов для эвакуации сотрудников из здания при
возникновении ЧС. Программа для ЭВМ.

{\bfseries References}

1. Kamenskaja E.N. Chrezvychajnye situacii mirnogo i voennogo vremeni. -
Rostov-na-Donu, Taganrog: Izd-vo Juzhnogo FU, 2020. - 160 s. ISBN
978-5-9275-3489-0. {[}in Russian{]}

2. Gasnikov A.V. i dr. Vvedenie v matematicheskoe modelirovanie
transportnyh potokov. - M.: MCNMO, 2013. - 427 s. ISBN
978-5-4439-0040-7. {[}in Russian{]}

3. Birjukov V.V. Passazhirskie perevozki v gorodah i aglomeracijah. -
Novosibirsk: NGTU, 2020. - 368 s. ISBN 978-5-7782-4264-7. {[}in
Russian{]}

4. Safronov Je.A. Transportnye sistemy gorodov i regionov. - M.: Izd-vo
ASV, 2007. -- 288 s. ISBN 978-5-93093-345-1. {[}in Russian{]}

5. Pashkov N.N. Transportnaja logistika (linejnoe programmirovanie). -
M.: Prometej, 2020. - 202 s. ISBN 978-5-00172-021-8. {[}in Russian{]}

6. Zjabirov H.Sh., Shapkin I.N. Optimizacija, issledovanie operacij i
teorija upravlenija transportnymi processami. - M.: Finansy i
statistika, 2023. - 520 s. ISBN 978-5-00184-099-2. {[}in Russian{]}

7. Zjabirov H.Sh., Shapkin I.N., Jusipov R.A. Jeffektivnye metody i
modeli prognozirovanija transportnyh processov i sistem. - M.: Finansy i
statistika, 2024. - 298 s. ISBN 978-5-00184-114-2. {[}in Russian{]}

8. Mazakov T.Zh., Burgegulov A.D., Dzhomartova Sh.A., Mazakova A.T.,
Sametova A.A., Dosanalieva A.T. Postroenie marshrutov dlja jevakuacii
sotrudnikov iz zdanija pri vozniknovenii chrezvychajnoj situacii//
Vestnik KazUTB. - Astana, 2024. - № 1(22). - S.68-74. DOI
10.58805/kazutb.v.1.22-292. {[}in Russian{]}

9. Burgegulov A.D., Dzhomartova Sh.A., Mazakova A.T., Aliaskar M.S.,
Mazakov T.Zh., Majlybaeva A.D. Postroenie putej jevakuacii ljudej iz
mnogojetazhnogo zdanija pri ugroze i vozniknovenii pozhara// Vestnik
KazUTB. - Astana, 2025. - № 1(26). - C.126-133. DOI
10.58805/kazutb.v.1.26-903. {[}in Russian{]}

10. A.S. №42850 ot 13.02.2024 g. Mazakov T.Zh., Burgegulov A.D.,
Dzhomartova Sh.A., Tojkenov G.Ch., Majlybaeva A.D., Zijatbekova G.Z.,
Mazaқova Ә.T., Әliasқar M.S., Sametova A.A., Dosanalieva A.T. Programma
postroenija marshrutov dlja jevakuacii sotrudnikov iz zdanija pri
vozniknovenii ChS. Programma dlja JeVM. {[}in Russian{]}

\emph{{\bfseries Сведения об авторах}}

Демеубаева Д.К. - старший преподаватель МИТУ, Алматы, Казахстан, e-mail:
laura.demeubayeva1618@gmail.com;

Мазаков Т.Ж.- д.ф.-м.н., профессор МИТУ, Алматы, Казахстан, e-mail:
tmazakov@mail.ru;

Джомартова Ш.А.- д.т.н., профессор КазНУ им. аль-Фараби,
Алматы,Казахстан, e-mail:
jomartova@mail.ru;

Мазакова А.Т.- PhD, старший преподаватель КазНУ имени аль-Фараби,
Алматы, Казахстан, e-mail:
aigerym97@mail.ru;

Майлыбаева А.Д.- к.ф.-м.н., доцент, Атырауский университет им. Х.
Досмухамедова, Атырау, Казахстан, e-mail:
a.maylibayeva@asu.edu.kz;

Әлиасқар М.С.- старший преподаватель МИТУ, Алматы, Казахстан, e-mail:
m.alyasqar@gmail.ru;

Чжао Байдун - старший преподаватель, Даляньский политехнический
университет, Далянь, Китай., e-mail: xiaoyuezhishang@vip.qq.com;

Досаналиева А.Т.- сеньор-лектор, Алматинский Технологический
Университет, Алматы, Казахстан, e-mail: Dosanalieva1985@gmail.com.

\emph{{\bfseries Information about the authors}}

Demeubayeva L.- senior lecturer at IETU, Almaty, Kazakhstan, e-mail:
laura.demeubayeva1618@gmail.com;

Mazakov T.Zh. - Doctor of Physical and Mathematical Sciences, Professor
IETU, Almaty, Kazakhstan, e-mail: tmazakov@mail.ru;

Jomartova Sh.A.- Doctor of Technical Sciences, Professor, Al-Farabi
Kazakh National University, Almaty, Kazakhstan, e-mail:
jomartova@mail.ru;

Mazakova A.T.- PhD, senior lecturer, al-Farabi Kazakh National
University, Almaty, Kazakhstan, e-mail: aigerym97@mail.ru;

Mailybayeva A.- PhD, Associate Professor, Kh. Dosmukhambetov Atyrau
University, Atyrau, Kazakhstan, e-mail:
a.maylibayeva@asu.edu.kz;

Aliaskar M.-senior lecturer at IETU, Almaty, Kazakhstan, e-mail:
m.alyasqar@gmail.ru;

Zhao Baidong \emph{-} Senior Lecturer, Dalian Polytechnic University,
Dalian, China, e-mail: xiaoyuezhishang@vip.qq.com;

Dossanaliyeva A.- Senior Lecturer, Almaty Technological University,
Almaty, Kazakhstan, e-mail: dosanalieva1985@gmail.com.

{\bfseries \hfill\break
}\
